\section{引言}\label{sec:intro}

条件概率递推估计（原称 Bayesian estimation）为动态状态估计问题提供了严格的一般框架。该方法的基本思路，是基于截至当前时刻的全部可用信息，构造待估状态的条件概率密度函数（PDF）。对于线性正态估计问题（原称 linear Gaussian estimation），所需的条件概率密度函数在滤波器的每一次迭代中均保持为正态分布（原称 Gaussian distribution），此时线性最小方差递推滤波器（原称 Kalman filter）的递推关系恰好精确传播并更新该分布的均值与协方差。然而，对于一般非线性或非正态问题，所需状态条件密度通常不存在解析（闭式）表达式。

一阶线性化正态滤波器（原称 extended Kalman filter）是递推非线性状态估计中最常用的方法 \cite{r01}。该方法在每一步围绕当前预测状态对非线性模型进行一阶微分线性化，从而借用线性最小方差递推滤波器的更新公式。在此框架下，所需的后验概率密度函数仍被强行近似为单峰正态分布；这一近似可能严重扭曲真实的底层分布结构，并导致滤波器发散。其他解析近似途径包括正态混合和滤波器（原称 Gaussian sum filter） \cite{r02}，以及仅近似分布前两阶矩的方法 \cite{r03,r04}。更为直接的数值求解思路是在状态空间中构造离散网格以数值计算所需概率密度函数 \cite{r05,r06,r07,r08,r09}。然而，设计高效网格本身颇为棘手，且在多维状态空间中往往需要极其庞大的网格点数，每个网格点上均须执行可观的计算量。

本文提出一种新方法，用以表示并递推生成状态条件概率密度函数的近似。其核心思想是将所需的概率密度函数表示为一组随机样本（粒子），而非状态空间上的显式函数 \cite{r10,r11,r12}。当样本规模足够大时，这组样本在统计意义上构成所需概率密度函数的一种精确、等价的经验表示。状态向量分布的各阶矩估计（如均值与协方差）以及分位数估计均可直接由样本算得。如有需要，还可由样本重建概率密度函数的连续函数估计 \cite{r13}，进而给出最高后验密度（HPD）区域或众数的估计。文献 \cite{r11} 曾提出一种自适应重要性抽样算法，但该方法会生成庞大的混合分布，且计算代价高昂。

本文提出并在第~\ref{sec:bootstrap}~节 中论证一种递推采样算法——称为条件自举粒子滤波器（下文简称条件自举滤波器或自举滤波器，原称 bootstrap filter），用于针对第~\ref{sec:bayes}~节 所述的离散时间估计问题向前传播并更新样本集。该滤波器之所以如此命名，是因为算法最关键的后验更新阶段（即条件后验更新法则，原称 Bayes' rule）是通过加权自举重采样实现的，其思路受到 Smith 与 Gelfand \cite{r14} 所证结果的启发；Müller \cite{r12} 亦描述过一种相关算法。这类基于随机抽样的技术彻底免除了在状态空间中预先划分网格的需要：样本会自然地向高概率密度区域聚集。此类方法的另一显著优势在于，能够直接处理任意形式的非线性函数以及任意分布类型的系统过程噪声与观测噪声。

第~\ref{sec:sim}~节 给出了两个数值仿真算例，其中第二个算例为四维状态空间下的纯方位目标追踪问题。文中系统比较了条件自举滤波器与标准一阶线性化正态滤波器的估计性能。
